\documentclass[a4j, 9pt,uplatex,dvipdfmx]{jsarticle}
\usepackage{color}
\usepackage[hiresbb]{graphicx}
% \usepackage[utf8]{inputenc}
% \usepackage{ascmac}
\usepackage{amsfonts}
\usepackage{amsmath}
\usepackage{amssymb}
\usepackage{amsthm}
\usepackage{comment}
\usepackage[shortlabels]{enumitem}
\usepackage{here}
\usepackage{latexsym}
\usepackage{multicol}
\usepackage{url}
\usepackage[hidelinks]{hyperref}
\usepackage{pxjahyper}
\usepackage[style=ieee, citestyle=numeric, backref=false, backend=biber]{biblatex}
\addbibresource{ref.bib}

\usepackage{sakolab-summary}

% 表紙の項目
% \summarytype{修 士 論 文 概 要 書}
% \esummarytype{Summary of Master's Thesis}
\date{07/20/2026} % MM/DD/YYYYの形式で記入
\department{情報理工・\\\fontsize{10truept}{10truept}\selectfont 情報通信専攻}
\author{Takumi Otsuka}
\advisor{Kazue Sako}
\guidance{暗号プロトコル\\\fontsize{10truept}{10truept}\selectfont 研究}
\studentid{5124FG15}
\title{Locating the Deployability Wall for Post-Quantum Anonymous Credentials in General-Purpose zkVMs}
\authorsize{12truept}
\titlesize{8.5truept}

\begin{document}
    % 二段組(タイトルを除く)
    \maketitle

    \begin{multicols*}
        {2}
        % 本文
		% NOTE (yours to reframe): this is a fact-correct scaffold, not final framing.
% The goal sentence and motivation are where your taste sits -- rewrite in your voice.
\section{Introduction}
Post-quantum anonymous credentials let a holder prove possession of a valid
credential without revealing identity, under assumptions that survive a
quantum adversary. We target the setting in which the holder generates this
proof on a \emph{personal device}, here a consumer laptop (18-core, 24\,GB),
as opposed to a cloud or server prover. A general-purpose zero-knowledge
virtual machine (zkVM) is attractive in this setting for \emph{flexibility}
rather than for raw speed: it proves an arbitrary verification predicate without
a bespoke, per-scheme circuit. The obstacle is arithmetic. The schemes we
care about (PLUM, Loquat) operate over a $192$-bit prime field, while
production zkVMs (RISC~Zero, SP1) fix a small ($\approx 31$-bit) native
field, so every operation is emulated by multi-limb arithmetic; an algebraic
hash (Griffin) accounting for $\approx 91\%$ of the verification cost
amplifies the penalty. This thesis asks whether such a verification can be
made to run on a personal device at all, what it costs, and whether the
resulting proof can additionally be zero-knowledge.

        % \input{sections/02-preliminaies}
        \section{The Field-Mismatch Tax and When a Precompile Pays}
We formalize the \emph{field-mismatch tax} as the performance gap that arises when a cryptographic primitive's natural algebraic domain (e.g., arithmetic over a large prime field $\mathbb{F}_q$ or structured rings) does not coincide with the proof system's \emph{native} arithmetic, and derive from it a criterion for when a custom precompile lowers proving cost. We model a proof backend by its \emph{arithmetization regime}: (i) a native field $\mathbb{F}_p$, (ii) an encoding of the computation domain into $\mathbb{F}_p$ (possibly via multi-limb representations), and (iii) a constraint/execution representation (e.g., R1CS for static circuits or AIR/trace constraints for STARK-based zkVMs). This viewpoint separates two qualitatively different regimes.

In the \emph{field-aligned} regime typical of application-specific static circuits, the designer can select (or align to) a proof system whose scalar field matches the primitive's domain, so arithmetic operations are expressed directly: multiplications correspond to constant-cost constraints, and total constraint growth tracks the primitive's multiplicative complexity. In the \emph{mismatch} regime typical of general-purpose zkVMs, the native field $\mathbb{F}_{p_{\text{VM}}}$ is fixed by the VM/STARK design and is often much smaller than $\mathbb{F}_q$. Elements of $\mathbb{F}_q$ must then be represented as $k=\lceil \log_2(q)/w\rceil$ $w$-bit limbs, requiring carry propagation and modular reduction to emulate non-native arithmetic. We show that even a single non-native modular multiplication is lower-bounded by $\Omega(k)$ native multiplications (with schoolbook emulation realizing all $k^2$ cross-products, plus carries and reduction), yielding an inherent field-mismatch tax that cannot be eliminated by constant-factor engineering alone. The lower bound is a counting argument: in a limb-aligned routine the product depends on every one of the $k$ input limbs, while each native multiplication carries at most one of them, so at least $k$ native multiplications are forced.

For STARK-based zkVMs, this cost further amplifies at the proof-system level: limb-based arithmetic increases the executed instruction count and inflates the AIR execution trace length by $\Theta(k^2)$ under the schoolbook emulation the evaluated zkVMs use. Since zkVM prover time is dominated by low-degree testing and commitments and scales roughly as $O(W\cdot T\log T)$ in trace width $W$ and length $T$, arithmetic mismatch translates into a fundamental prover slowdown relative to field-aligned static circuits. This yields the thesis's criterion: a precompile lowers proving cost only on the field-mismatch axis, and it pays only when it removes more emulated trace area than the dedicated chip itself adds. A primitive already served by an existing precompile, or one native to the prover field, therefore does not warrant a bespoke chip even though one could be built.

        \section{System Design}
We construct a custom precompile around one reusable interface and instantiate
it as a small \emph{family} of three chips over PLUM's $199$-bit prime field:
(i) a bespoke Griffin permutation~\cite{cryptoeprint:2022/403} AIR, the
load-bearing $91\%$ target; (ii) a dedicated \texttt{FP192\_MUL}
field-multiplication chip; and (iii) a \texttt{FP192\_POW\_RES} chip composing
that multiplication into the $t$-th power-residue PRF symbol check
$a^{(p-1)/t}\bmod p$ ($t=256$ fixed). All three build and pass a core-STARK
smoke proof (execute, prove, verify), each with a written soundness argument
tying its constraints to the reference specification, so the precompile layer
preserves the scheme's unforgeability under the assumption that the zkVM's
cross-table lookup argument is sound. Only the Griffin chip is on the measured
path; the field-multiplication chip does not pay (the host's $256$-bit
modular-multiplication precompile already serves it), while whether the
power-residue chip---the one irreducibly-large-prime primitive---pays is an open
question the criterion poses precisely.

        \section{Evaluation}
We measure four cells on the target hardware: PLUM verification in a zkVM
(1)~without and (2)~with the Griffin precompile; (3)~the same verification
using software SHA-3 as a control; and (4)~PLUM verification in an
application-specific circuit (Aurora) as an in-SNARK reference.

Without the precompile, verification does not terminate within a
$24\,\mathrm{GB}/3\,\mathrm{h}$ bound ($\approx 7.08$ billion cycles at
$\lambda=80$, projecting to tens of hours). With the precompile it completes
in roughly \fbox{XX}~min on the $24\,\mathrm{GB}$ machine at a peak
$18.67\,\mathrm{GB}$ resident---the first configuration in the study to fit
the budget.
% DECISION: \fbox{XX} time is unsettled -- 32.5 min (2026-05-20) vs 14.89 min
% (2026-06-18, re-measured clean). 14.89 is more consistent with the 1.13x cycle
% ratio vs Cell 3's 13.3 min, but cause unverified; thesis prose still cites 32.5.
% Reconcile before submission, then replace this placeholder.

The control is the sharpest result. Even with the precompile, PLUM
verification executes $1.13\times$ more cycles ($125.5$\,M) than the same
verification using software SHA-3 ($110.8$\,M): a standardized hash,
unaccelerated, still beats our accelerated algebraic hash. This is the
\emph{precompile paradox} made empirical---the cost of choosing a
SNARK-friendly custom primitive over one the ecosystem has standardized.

        \section{Discussion}
Feasibility is not deployability. A succinct proof on this substrate is not
zero-knowledge: SP1's default STARK proof is succinct but not ZK, and ZK is
obtained only by wrapping it in a Groth16/PLONK SNARK over a pairing-friendly
curve---which is \emph{not} post-quantum, so a future quantum adversary
retroactively de-anonymizes every stored transcript. The post-quantum ZK
route (a masked zk-STARK) yields only \emph{computational} post-quantum
anonymity in the (Q)ROM, not everlasting (information-theoretic) anonymity,
because the witness trace is committed under a computationally-hiding hash.
No available zkVM path delivers post-quantum soundness and zero-knowledge
jointly. This \emph{deployability wall}---not the prover's speed---is the
binding obstacle for a personal-device, post-quantum anonymous credential.

        \section{Conclusion}
We formalized the arithmetic mismatch as a quantifiable agility tax,
specified a precompile suite that brings post-quantum anonymous-credential
verification within the personal-device memory budget on this substrate, and
showed both an empirical instance of the precompile paradox and a
deployability wall that blocks a jointly post-quantum and zero-knowledge
proof. Closing that wall (an everlasting-anonymous proof on a transparent
post-quantum substrate, and pushing below the personal-device budget toward
mobile) is left as future work.


    	% \section*{謝辞}
    	% 本研究の成果の一部は，公益財団法人セコム科学技術振興財団の特定領域研究「権限の濫用を抑止する技術的・社会的しくみに関する研究」によって得られたものです．ここに感謝の意を表します．


	    %
	    %  参考文献
	    %  ref.bibを用意する。
	    %
	    % \bibliographystyle{junsrt}
	    % \bibliography{ref}
		\printbibliography[title=参考文献]
	\end{multicols*}
\end{document}